\begingroup%
\definecolor{ksBlue}{HTML}{3478B8}%
\definecolor{ksTeal}{HTML}{27877F}%
\definecolor{ksAmber}{HTML}{C58B28}%
\definecolor{ksInk}{HTML}{243447}%
\definecolor{ksPaleBlue}{HTML}{F0F6FC}%
\definecolor{ksPaleTeal}{HTML}{EFF8F5}%
\definecolor{ksPaleAmber}{HTML}{FFF8E8}%
\begin{tikzpicture}[
  x=1mm, y=1mm,
  card/.style={draw, fill=white, rounded corners=2mm, line width=.7pt,
    text=ksInk, align=flush center, font=\fontsize{9.5}{12.5}\selectfont,
    text width=69mm, minimum width=73mm, minimum height=16mm, inner xsep=2mm, inner ysep=1.5mm,
    outer sep=.35pt},
  challenge/.style={card, font=\bfseries\fontsize{9.5}{12.5}\selectfont},
  frame/.style={draw, rounded corners=3mm, line width=.7pt, inner sep=3mm,
    outer sep=.35pt},
  heading/.style={text=ksInk, font=\bfseries\fontsize{10}{12.5}\selectfont,
    inner sep=0pt, outer sep=0pt},
  link/.style={line width=.65pt},
  flow/.style={-{Stealth[length=2.6mm,width=1.9mm]}, ksInk, line width=.9pt},
  edge text/.style={text=ksInk, font=\fontsize{9}{12.5}\selectfont,
    inner sep=1mm}
]
  \node[heading] (s1-title) at (40,0) {阶段一：事前 · 数据准备与表示层};
  \node[challenge, draw=ksBlue] (c1-title) at (0,-13)
    {挑战一：硬件感知弱、编程能力强失衡};
  \node[card, draw=ksBlue] (c1-a) at (80,-13)
    {语料偏置混杂：NV 历史架构主导且跨代混杂\\
    新架构 (Blackwell) 与国产异构卡严重匮乏};
  \node[card, draw=ksBlue] (c1-b) at (80,-32)
    {语言杂（语言即思维边界）：\\
    CUDA / CuTeDSL / Triton /\\
    TileLang / Ascend / HIP 各自割裂};
  \node[card, draw=ksBlue] (c1-c) at (0,-32)
    {内容杂（数据未压缩未蒸馏）：\\
    Host / IO 样板淹没核心设计，\\
    高熵源码阻碍软硬件协同规律涌现};

  \node[heading] (s2-title) at (40,-56) {阶段二：事中 · 生成与迭代过程层};
  \node[challenge, draw=ksTeal] (c2-title) at (0,-69)
    {挑战二：迭代收敛难、最优点与停点不明};
  \node[card, draw=ksTeal] (c2-a) at (80,-69)
    {高随机性与低一致性：\\
    完全相同条件下重复运行轨迹剧烈分化，\\
    性能帕累托前沿模糊};
  \node[card, draw=ksTeal] (c2-b) at (80,-88)
    {设计决策 + 代码实现强耦合：\\
    编译报错或机群波动即引发路径分叉，\\
    好流水因局部实现瑕疵被误杀};
  \node[card, draw=ksTeal] (c2-c) at (0,-88)
    {避难就易走弯路：为求快速跑通\\
    退回 mma.sync 等过时指令，平台期漫长\\
    且不知何时停止投入预算};

  \node[heading] (s3-title) at (40,-112) {阶段三：事后 · 结果对象与交付维护层};
  \node[challenge, draw=ksAmber] (c3-title) at (0,-125)
    {挑战三：介入理解难、生产维护成本高};
  \node[card, draw=ksAmber] (c3-a) at (80,-125)
    {Shape 分支特化膨胀：靠大量 if-else 派生拟合\\
    单一负载而非泛化设计突破，代码急剧膨胀};
  \node[card, draw=ksAmber] (c3-b) at (80,-144)
    {隐式缺陷暗藏：共享内存重复写入竞争、\\
    边界越界与死锁隐患频发，\\
    迫使测试集膨胀、迭代变慢};
  \node[card, draw=ksAmber] (c3-c) at (0,-144)
    {纯黑盒阻断专家介入：\\
    高层设计取舍被底层代码抹平\\
    人工难审阅、上线风险高，实用性受阻};

  \begin{scope}[on background layer]
    \node[frame, draw=ksBlue, fill=ksPaleBlue,
      fit=(s1-title)(c1-title)(c1-a)(c1-b)(c1-c)] (s1) {};
    \node[frame, draw=ksTeal, fill=ksPaleTeal,
      fit=(s2-title)(c2-title)(c2-a)(c2-b)(c2-c)] (s2) {};
    \node[frame, draw=ksAmber, fill=ksPaleAmber,
      fit=(s3-title)(c3-title)(c3-a)(c3-b)(c3-c)] (s3) {};
  \end{scope}

  \draw[link, ksBlue] (c1-title) -- (c1-a);
  \draw[link, ksBlue] (c1-a) -- (c1-b);
  \draw[link, ksBlue] (c1-b) -- (c1-c);
  \draw[link, ksTeal] (c2-title) -- (c2-a);
  \draw[link, ksTeal] (c2-a) -- (c2-b);
  \draw[link, ksTeal] (c2-b) -- (c2-c);
  \draw[link, ksAmber] (c3-title) -- (c3-a);
  \draw[link, ksAmber] (c3-a) -- (c3-b);
  \draw[link, ksAmber] (c3-b) -- (c3-c);
  \draw[flow] (s1.south) --
    node[edge text, right=1mm] {带着残缺先验进入搜索} (s2.north);
  \draw[flow] (s2.south) --
    node[edge text, right=1mm] {产出高风险黑盒代码} (s3.north);
\end{tikzpicture}%
\endgroup%
